\section{引言}

燃放烟花爆竹会在一夜之间把一座城市的 \pmm{2.5} 推高数倍。这件事在公共财政学的框架里有两点值得注意。

第一，它是\textbf{消费}产生的负外部性，不是生产。清洁空气是非排他、非竞争的公共品，燃放者不会把对同城居民的损害计入自己的决策，因此市场供给必然不足，	extbf{政府干预有明确的效率理由}。这也排除了绝大多数环境政策工具惯用的着力点——你无法对烟囱加装除尘器，也无法要求消费者“清洁燃烧”；唯一的减排方式是不烧。

第二，它的损害在时间与空间上都极端集中。集中于一年之中的十几个小时，集中于城市内部的居民区。绝大多数环境外部性是弥散的、连续的、按年计量的；这一个不是。

中国对这件事有两套并行的制度安排。

一套是\textbf{价格工具}。消费税税目“鞭炮、焰火”按 15\% 从价征收，在生产、委托加工与进口环节缴纳。该税目自 1994 年 1 月 1 日消费税开征时即已存在，是 1993 年《消费税暂行条例》原始 11 个税目之一。2006 年 4 月那次把消费税明确转向资源节约与环境保护的调整，新增并给出环保理由的是木制一次性筷子与实木地板，\textbf{没有触及该税目}。1994 年该税目的税率同样是 15\%。也就是说，\textbf{这个带着环境理由设立的价格工具，三十余年没有被重新校准过一次}。

另一套是\textbf{数量工具}。自 1990 年代起，各城市陆续出台禁放、限放规定；2017 年国家大气污染防治方案曾明确要求各城市制定禁放限放方案，2023 年的国家方案则删去了这一条，城市层面的管控转而挂靠在地方自行印发的大气污染防治文件中。34 个城市的文件取证显示，2023 年之后城市写不写这一条，是地方自主加码的结果，而不是上级的硬性要求（见第~\ref{sec:institution}~节）。

这不是一个“某项政策有没有效”的问题，而是一个关于\textbf{工具形式}的一般性问题：当损害的变异集中在城市与时间两个维度、而工具的设定层级与适用维度与之不匹配时，会出现什么后果。烟花爆竹只是一个恰好把这一矛盾放大到可见程度的场景。

于是本文的问题可以写成一句话：\textbf{在已经存在一个价格工具的情况下，为什么地方政府持续叠加数量工具，而那个价格工具却三十年未被调整？}

本文的回答分三步。

\textbf{第一步是测量外部性本身。} 用 389 个城市、2015—2026 共 12 个春节窗口的逐小时数据，本文刻画出尖峰的量级、位置、城际异质性与污染物结构，并进一步把尖峰增量在城市与年份两个维度上做方差分解。这一步不依赖任何政策变异，因此不受后文识别失败的影响，构成本文最硬的部分。

\textbf{第二步是尝试因果识别，并报告它的失败。} 本文用交错双重差分估计禁放对初一日均 \pmm{2.5} 的影响。结论是这个设计\textbf{拿不出可信的估计}。失败的证据是递进的：可用作识别的市级切换城市只有 6 个已核实（另有一个未核实），且三类切换方向不一；89\% 的城市-年没有政策记录；安慰剂结果变量显著；安慰剂“日”扫描显示估计量对燃放日没有亲和性；而根因在于禁令总是与同期的工业减排方案同文下达，二者在数据上不可分离。

\textbf{第三步是回到工具本身。} 本文主张，问题的关键不在税率高低，而在\textbf{工具的解析度}。一个全国统一、按年适用、在生产环节征收的从价税，无论税率定得多高，都无法区分“哪个城市”“哪一天”；而第~\ref{sec:magnitude}~节的测量恰恰显示，损害的大部分变异落在城市 $\times$ 年份这一维度上。\textbf{提高税率改变的是价格水平，不改变工具的解析度。}

这一步也构成对既有文献的正面回应。中文财税文献对“鞭炮、焰火”税目并不陌生：安体富、樊勇（2006）把它列入“过度消费危害人类健康、社会秩序和生态环境的特殊消费品”，并建议调高税率；庞凤喜、王绿荫（2020）明确写“鞭炮焰火……消费税税率偏低”，主张提高税率并后移征税环节。本文不否认税率偏低，但与这一主张的\textbf{方向}不同：如果约束条件是解析度而非价格水平，那么提高税率改变的是价格水平，不改变解析度；它仍可能是一个粗糙的平均抑制信号，但不能替代解析度改革。

本文的贡献因此不是“第一次研究这个问题”——既有文献已经用全国面板做过春节污染（Li et al., 2025；Chen et al., 2022），也已经多次点名这个税目。本文的贡献是三点：把外部性的\textbf{时间—空间结构}测准，并给出一个可复算的方差分解；如实报告一次\textbf{失败的因果识别}及其结构性原因；把“工具成本的可见性”这一已有框架（Beiser-McGrath et al., 2023；Harrington et al., 2004；Oates \& Baumol, 1975）应用到一个新的工具对与新的场景上。
